% ---------------------------------------------------------------------------
% The core of this version. Theorem labels are the complete paper's, because the
% four proof blocks under sections/moved/ are input verbatim: one copy of each
% proof, so the two documents cannot drift. Numbers come from notebook 04 and
% data/processed/regex_glob.csv.
% ---------------------------------------------------------------------------

\section{The Regular Sub-problem}
\label{s:sec:regular}

Both languages filter on tag values, and both do it with a pattern language, but not the
same one. Datadog's is a glob; PromQL's is a regular expression. One is a proper subclass
of the other, so the two directions of this sub-problem are not the same question asked
twice: the forward direction needs a translation that always exists, and the reverse
direction needs a decision procedure for whether one does. This section gives both, and
the decision procedure is where the course's machinery does the work.

\subsection{Globs, and what they can say}
\label{s:sec:globs}

\begin{definition}[Tag glob]\label{def:glob}
Fix an alphabet $\Sigma$ of tag-value characters. A \emph{glob} is a string over
$\Sigma\cup\{*\}$, written $g=s_0*s_1*\cdots*s_n$ with $n\ge 0$ stars and literal
segments $s_i\in\Sigma^{*}$. It denotes
\[
L(g)=\{\,s_0w_1s_1w_2\cdots w_ns_n \;:\; w_i\in\Sigma^{*}\,\}.
\]
Datadog documents prefix, suffix and substring wildcards and no other
construct~\citep{ddfilter}: there is no single-character wildcard, no alternation inside
a value, no character class, and no escape that would make a literal \texttt{*}. That
sentence is the whole premise of this section's negative results, and it was read from
the documentation until we put it to a real account instead --- the account was asked for
all three constructs and the query parser rejected all three, which is a stronger answer
than selecting nothing. Write $m(g)=\sum_i|s_i|$ for the total literal length.

Which set $\Sigma$ is matters, and is fixed in \S\ref{s:sec:glob-corpus} rather than
assumed here. Every language in this section is read over it, including the $\Sigma^{*}$
that Proposition~\ref{prop:glob-regex} reads \texttt{.*} as, and the two available
readings differ on $238$ of this corpus's matcher occurrences.
\end{definition}

The forward direction needs only the easy half, and gets it by construction.

\begin{proposition}[Globs are regular, uniformly]\label{prop:glob-regex}
Let $\mathrm{glob}(g)$ be the regular expression obtained from $g$ by replacing every
\texttt{*} with \texttt{.*} and escaping every other character. Then
$L(\mathrm{glob}(g))=L(g)$, where the regular expression is read under PromQL's
whole-value anchoring.
\end{proposition}

\begin{proof}[Proof of Proposition~\ref{prop:glob-regex}]
Induction on $n$. For $n=0$ both sides are $\{s_0\}$. For the step, \texttt{.*} denotes
$\Sigma^{*}$ and concatenation of regular expressions denotes concatenation of languages,
which is exactly how Definition~\ref{def:glob} quantifies $w_i$. Anchoring matters: it is
what makes the outermost concatenation span the whole value, and it is the fact
Proposition~\ref{prop:disjunction} rests on, read off a real server there rather than
assumed.
\end{proof}


\noindent
Every glob therefore has an exact PromQL matcher, which is why no gap code in
Table~\ref{s:tab:gaps} mentions globs on tag \emph{values}. The remaining direction is
the interesting one, and one structural fact decides most of it.

\begin{lemma}[Dichotomy]\label{lem:dichotomy}
For every glob $g$ over a non-empty $\Sigma$: if $n=0$ then $L(g)=\{s_0\}$ is a single
word; and if $n\ge 1$ then $L(g)$ contains a word of every length $\ell\ge m(g)$, and
every word of $L(g)$ contains each $s_i$ as a factor, in order.
\end{lemma}

\begin{proof}[Proof of Lemma~\ref{lem:dichotomy}]
The first case is immediate. For the second, taking $w_1\in\Sigma^{\ell-m(g)}$ and
$w_i=\varepsilon$ for $i>1$ gives a word of length $\ell$; the factor claim is
Definition~\ref{def:glob} read left to right.
\end{proof}


\noindent
So a glob language is either a singleton or unbounded in length. That is a strong
restriction, and together with what a star admits --- every character of $\Sigma$, by
Definition~\ref{def:glob} --- it settles each construct the corpus refuses, not by
measurement but by proof. The second ingredient is easy to lose sight of: parts~(1)
and~(2) below come from the dichotomy alone, and part~(3) does not.

\begin{corollary}[Three things no glob says]\label{cor:no-glob}
Assume $|\Sigma|\ge 2$. Then
\begin{enumerate}
\item no glob denotes a finite language with two or more words --- in particular none
denotes $\Sigma$, so there is no single-character wildcard, and none denotes an
alternation of distinct literals;
\item no glob denotes $\Sigma^{+}$, the non-empty values;
\item for $\emptyset\neq A\subsetneq\Sigma$, no glob denotes $A^{+}$ --- a star cannot be
restricted to part of the alphabet.
\end{enumerate}
\end{corollary}

\begin{proof}[Proof of Corollary~\ref{cor:no-glob}]
(1) A finite language with two words is neither a singleton nor unbounded, so
Lemma~\ref{lem:dichotomy} excludes it. (2) $\varepsilon\notin\Sigma^{+}$ forces
$m(g)\ge 1$, since $m(g)=0$ makes $g$ a string of stars and $\varepsilon\in L(g)$; and
$\Sigma^{+}$ contains words of length $1$, so $m(g)\le 1$. Hence $m(g)=1$: $g$ has a
single literal character $c$ and at least one star. By Lemma~\ref{lem:dichotomy} every
word of $L(g)$ contains $c$, while $\Sigma^{+}$ contains the one-character words of
$\Sigma\setminus\{c\}$, which is non-empty. (3) The same argument gives $m(g)=1$ with one
literal $c$, which must lie in $A$ since every word of $L(g)$ contains it. A single
literal and at least one star leaves three shapes: $c{*}$, ${*}c$ and ${*}c{*}$. In the
first and third $L(g)$ contains $c\sigma$ for every $\sigma\in\Sigma$; in the second it
contains $\sigma c$ for every $\sigma$. Either way it contains a word using a
$\sigma\notin A$, which $A^{+}$ does not, and $\Sigma\setminus A\neq\emptyset$.

Taking $c\sigma$ in every case, as an earlier version of this proof did, is wrong for
$g={*}c$: that language has no word beginning with $c$ unless $c$ repeats.
\end{proof}


\noindent
Part~(3) generalises, and the generalisation is what the unconditional refusals of
\S\ref{s:sec:glob-corpus} rest on, so it is stated rather than left as an aside.

\begin{corollary}[A character a language never uses]\label{cor:absent-char}
Let $L$ be infinite and let some $\sigma\in\Sigma$ occur in no word of $L$. Then $L$ is
not a union of finitely many glob languages.
\end{corollary}

\begin{proof}[Proof of Corollary~\ref{cor:absent-char}]
A finite union covering an infinite $L$ has a member with infinite language, which by
Lemma~\ref{lem:dichotomy} is not a singleton and so contains a star. A star admits every
character of $\Sigma$, so that member's language contains a word using $\sigma$, and
that word lies in $L$ --- contradicting the assumption that $\sigma$ occurs in no word
of $L$.
\end{proof}


\noindent
This is the one result here that holds against an \emph{arbitrary} union, with no bound
on how many branches a translator may spend, which is why \S\ref{s:sec:glob-corpus} can
separate ten refusals from the rest.

\subsection{Deciding whether a regular expression has a glob}
\label{s:sec:decide}

Corollary~\ref{cor:no-glob} rules particular shapes out. The general question --- given a
regular expression, is there \emph{any} glob denoting the same language? --- is settled by
a single observation about shortest words.

\begin{proposition}[Glob-expressibility is decidable]\label{prop:decidable}
Let $r$ be a regular expression with $L(r)\neq\emptyset$, read \texttt{*} as matching the
empty string, and read $L(r)$, the candidate globs and the automata below over one fixed
non-empty alphabet $\Sigma$ --- the same $\Sigma$ Definition~\ref{def:glob} defers and
\S\ref{s:sec:glob-corpus} fixes. It is decidable whether some glob $g$ has $L(g)=L(r)$,
and a witness is produced when one exists.
\end{proposition}
\begin{proof}
Suppose $L(r)=L(g)$ for $g=s_0*\cdots*s_n$. Taking every $w_i=\varepsilon$ gives the word
$s_0s_1\cdots s_n\in L(g)$ of length $m(g)$, and every other word of $L(g)$ has some
$w_i\neq\varepsilon$ and so is strictly longer. Hence $s_0s_1\cdots s_n$ is the shortest
word of $L(g)$ \emph{and the only word of that length}.

So the procedure is: construct a DFA for $r$ --- Rich's regular-expression-to-FSM
construction (Thm 6.1) followed by the subset construction (Thm 5.3) --- and find its
shortest accepted word $w$ by breadth-first search from the start state, together with
the number of accepting paths of that length. If $L(r)$ has no unique shortest word, no
glob exists. If it does, then $w=s_0s_1\cdots s_n$, so $g$ is $w$ with stars inserted at
some subset of its $|w|+1$ gaps: at most $2^{|w|+1}$ candidates. Each is settled by
deciding $L(r)=L(\mathrm{glob}(g))$, which is decidable by Rich Thm 9.5. The search is
finite and exhaustive, so it answers either way.
\end{proof}

\noindent
The witness is what makes this useful rather than merely true: a decision procedure that
returned only a bit would tell a migration that a glob exists without telling it which
one. \S\ref{s:sec:glob-corpus} runs this procedure over the corpus and uses it to check
the refusals of the syntax-directed proposer rather than trust them, which is
Remark~\ref{s:rem:refusal}'s standard applied to this paper's own output.

\subsection{Unions, and the bound}
\label{s:sec:glob-unions}

A single glob is the wrong target. By Proposition~\ref{prop:disjunction} a Datadog scope
can disjoin over one tag key, and read in this direction that proposition says the reverse
translation of a regular expression should be sought as a \emph{finite union} of globs;
$\mathtt{LIST|GET}\mapsto k \in (\texttt{LIST},\texttt{GET})$ is then immediate.

The union must be bounded, and the bound is part of the claim rather than an
implementation detail. Every finite language is a union of singleton globs, one per word,
so boundedness is what stops the question from being trivial --- and it is not only a
matter of taste. By Corollary~\ref{cor:no-glob} no glob denotes $\Sigma^{+}$, yet
$\bigcup_{c\in\Sigma}L(c*)=\Sigma^{+}$, so the non-empty values \emph{are} a finite union
of globs: one for every character of the alphabet. A scope is a piece of text somebody
maintains, so a union of a million branches is not a translation.

We fix $B=|\Sigma|^{2}$ --- two unconstrained character positions --- and report anything
needing more as a refusal, naming the bound. What makes that the choice rather than a
round number is the corpus's own distribution. The union sizes this corpus actually needs
stop at $206$ and then jump straight to $|\Sigma|^{2}$, the signature of two free
positions; nothing in between occurs at all. So every bound from $206$ up to
$|\Sigma|^{2}-1$ admits exactly the same patterns, and the square is the next choice that
is a different decision rather than a different number. The bound is about maintenance and
not about the service: a real account was asked how long a scope it will accept, and the
answer is about four times what $B$ needs.

Whether glob-union-expressibility is decidable without a bound we leave open. The
shortest-word argument does not transfer: each glob of a minimal union contributes its own
shortest word, and nothing bounds those by the shortest word of the union, so the
candidate space is no longer finite. Within the bound the procedure used here is sound but
not known to be complete --- it proposes a union from the regular expression's syntax and
then decides equivalence, so what it returns is exact, while a refusal is in general a
statement about the proposer. Proposition~\ref{prop:decidable} removes that caveat for the
single-glob case.

\subsection{What the corpus contains}
\label{s:sec:glob-corpus}

\begin{table}[htbp]\centering\footnotesize
\begin{tabular}{@{}lrrrr@{}}
\toprule
Outcome & Distinct & Share & Occurrences & Share \\
\midrule
One glob & 61 & 33.2\% & 309 & 29.7\% \\
A finite union of globs & 73 & 39.7\% & 344 & 33.0\% \\
Neither ($\bot$) & 50 & 27.2\% & 388 & 37.3\% \\
\emph{Template-valued, not a regex} & 174 & --- & 7,546 & --- \\
\bottomrule
\end{tabular}

\caption{Every \texttt{=\textasciitilde{}} and \texttt{!\textasciitilde{}} matcher in the
PromQL half of the corpus. Shares are of the concrete regular expressions; the
template-valued row is counted but not classified, because the pattern is not in the query
text.}
\label{s:tab:regex-glob}
\end{table}

The first result is not about regular languages. Of $8{,}587$ regex matchers, $7{,}546$
--- $87.9\%$ --- have a Grafana template variable as their pattern, as in
\texttt{\{pod=\textasciitilde{}"\$pod"\}}. The value is chosen when somebody looks at the
dashboard, so there is no pattern to convert. For these the reverse direction is not hard
but \emph{statically undefined}, and that is a property of how dashboards are written
rather than of either language.

Among the $1{,}041$ occurrences whose pattern is concrete,
Table~\ref{s:tab:regex-glob} gives the three-way split. A glob exists for $957$ of them
($91.9\%$), but for more than half of those only as a union: one glob covers $309$
occurrences ($29.7\%$) and a bounded union covers $648$ ($62.2\%$). Treating a single glob
as the target --- the obvious reading of ``does this regex have a glob'' --- would
therefore have understated the reverse direction by the $62.2\%$ of occurrences a union
reaches and a single glob does not, and the construction that recovers them is
Proposition~\ref{prop:disjunction}, already proved for the other direction.

\begin{table}[htbp]\centering\footnotesize
\begin{tabular}{@{}lrrl@{}}
\toprule
Construct with no glob & Distinct & Occurrences & Shortest example \\
\midrule
No non-empty wildcard (\texttt{.+}) & 21 & 238 & \texttt{.+} \\
No single-character wildcard (\texttt{.}, \texttt{5..}) & 26 & 140 & \texttt{5..} \\
No bounded alphabet under a star (\texttt{[a-z]+}) & 3 & 10 & \texttt{5[0-9]+} \\
\bottomrule
\end{tabular}

\caption{The constructs behind the $84$ refused occurrences. The first is ruled out as a
single glob by Corollary~\ref{cor:no-glob} and the second by
Corollary~\ref{cor:absent-char}; only the second is also beyond a bounded union of globs.}
\label{s:tab:regex-glob-reasons}
\end{table}

\begin{table}[htbp]\centering\footnotesize
\begin{tabular}{@{}llr@{}}
\toprule
PromQL regex & Datadog scope value & Occurrences \\
\midrule
\texttt{\textasciicircum{}k8s\_.*} & \texttt{k8s\_*} & 60 \\
\texttt{.*external-dns.*} & \texttt{*external-dns*} & 36 \\
\texttt{.*pod.*} & \texttt{*pod*} & 24 \\
\texttt{LIST|GET} & \texttt{GET, LIST} & 41 \\
\texttt{ext.*|xfs} & \texttt{ext*, xfs} & 34 \\
\texttt{proxy|attach|log|exec|portforward} & \texttt{attach, exec, log, portforward, p\dots{}} & 32 \\
\texttt{.+} & $\bot$ No non-empty wildcard & 133 \\
\texttt{(/dev/)?(mmcblk.p.+|nvme.+|rbd.+|\dots{}} & $\bot$ No non-empty wildcard & 49 \\
\texttt{5..} & $\bot$ No single-character wildcard & 40 \\
\bottomrule
\end{tabular}

\caption{Worked cases, most frequent first: the PromQL pattern, the Datadog scope value
the reverse translation produces, and how many occurrences carry it. A $\bot$ row names
the reason rather than a scope. \texttt{LIST|GET} is the case of
\S\ref{s:sec:glob-corpus}'s closing paragraph on case folding: the scope is an exact
translation of the pattern and matches nothing in any Datadog database.}
\label{s:tab:regex-glob-examples}
\end{table}

\paragraph{Which alphabet.}
Every equivalence here is decided over an alphabet, and the question is not a formality:
``does this regular expression denote the same set as that glob'' has two readings, and on
this corpus they differ by $238$ matcher occurrences. Over strings from an unbounded
character set \texttt{.+} is no finite union of globs at all, so the $21$ patterns that
contain it are refused; over the strings a Datadog tag value may contain, \texttt{.+} is
the union of $c$\texttt{*} for each character $c$ --- $41$ branches --- and a scope can say
every one of them. Datadog documents that set --- lower-case letters, digits, underscores,
minuses, colons, periods and slashes, with the value lower-cased on ingest --- which is
$41$ characters. The corpus is censused rather than trusted: every distinct tag value in
the instantiated corpus, read through this paper's own Datadog parser, uses only those
$41$ characters, and every one of the $41$ is written somewhere in it. That is the right
universe, and not as a convenience: a string that is not a legal tag value cannot appear
on either side of the equivalence, so admitting it compares two languages on words neither
will ever be asked about.

That census is over the Datadog half of the corpus, and the patterns being decided come
from the PromQL half, which does not obey it. $34$ of the $184$ distinct patterns, $133$
of the $1{,}041$ concrete matcher occurrences or $12.8\%$, write literal characters
outside $\Sigma$: $33$ of them upper-case only --- \texttt{LIST|GET},
\texttt{POST|PUT|PATCH|DELETE}, \texttt{Pending|Running} --- and \texttt{<none>|Job}
adding angle brackets. What the decider substitutes $\Sigma$ for is the class \texttt{.}
denotes, and it leaves literals as written, so those $34$ are decided over $\Sigma$
together with their own literals. Both sides of each equivalence are read the same way, so
the verdicts stand as equivalences of languages and Table~\ref{s:tab:regex-glob}'s split
does not move.

The mismatch it hides is worth more than the bookkeeping. Datadog lower-cases a tag value
on ingest, so no value in any Datadog database is \texttt{GET}, and the scope
\texttt{k IN (GET, LIST)} that this section reports as an exact reverse translation of
\texttt{LIST|GET} matches nothing at all. Those $130$ upper-case occurrences are
expressible as globs and empty as queries. That is a second reverse-direction cost beside
the refusals of Table~\ref{s:tab:regex-glob-reasons}; it is not a fact about globs --- a
PromQL matcher on a mixed-case label value has no Datadog counterpart whatever the target
construct --- and this paper does not model it.

\paragraph{What refuses, and on what authority.}
With $\Sigma$ that set, Table~\ref{s:tab:regex-glob-reasons} has to be read against the
union target of \S\ref{s:sec:glob-unions} rather than against a single glob. Two
constructs account for every refusal, and neither is a single-character wildcard:
\texttt{5..} is the union of $|\Sigma|^{2}$ globs, one per pair of characters, and
\texttt{.+} the union of $|\Sigma|$ of them, and a scope can be made to say either. The
three impossibility results are untouched as theorems --- they are about a \emph{single}
glob, and Table~\ref{s:tab:regex-glob} is the measurement of how little of this corpus a
single glob reaches; what they no longer do, against a union target over a finite
alphabet, is refuse a corpus pattern on their own.

What refuses is each of them together with something else. Four patterns are proposed a
union of more branches than the bound admits: \texttt{2..|5..} is two free positions twice
over at $3{,}362$ branches against a bound of $41^{2}$, and the two that carry $71$ of
those $74$ occurrences are alternations over device names, each branch of which ends in
\texttt{.+} and costs $41$. Those refusals rest on a number this paper chose and states,
and they decide if the number is raised --- into a scope with thousands of branches in it,
which is the sense in which the bound is about maintenance rather than about language.
They also rest on the proposer: these are the branch counts of the union the
syntax-directed proposal produces, and nothing here establishes that no smaller union
denotes the same language, so ``$3{,}362$ branches'' is an upper bound on what the pattern
costs and not a lower bound on what it needs. \texttt{5[0-9]+} is different:
it would need a star that never admits a non-digit, forbidden by
Corollary~\ref{cor:absent-char} --- not by part~(3) of Corollary~\ref{cor:no-glob}, which
is about $A^{+}$ itself, and which \texttt{5[0-9]+} is not, its shortest word having three
characters where $A^{+}$ has words of one.

Only the second kind is unconditional: an infinite language needs a star, a glob's star
admits every character, and no union of finitely many such globs confines it to the
digits. So $74$ of the $84$ occurrences are refused by a theorem about single globs plus a
bound this paper chose, and $10$ by a theorem about globs outright.

\paragraph{The proposer checked against the decision procedure.}
For the single-glob question the same conclusion can be reached by measurement rather than
by proof. Re-deciding every concrete regular expression with the complete procedure of
Proposition~\ref{prop:decidable} reaches \emph{all} $184$ distinct patterns --- counts in
this paragraph are of patterns rather than of occurrences --- and disagrees with the
syntax-directed proposer on \emph{none} of them. It recovers the same glob for all $61$ of
the single-glob cases, and finds no glob for all $7$ refusals and all $116$ unions.
Nothing is left undecided: no pattern exceeds the search's own candidate bound, so the
agreement is over the whole population rather than over a reachable part of it, though
``complete'' is relative to the dialect the search parses.

\paragraph{One assumption the vendor does not document.}
Corollary~\ref{cor:no-glob}(2) assumed a glob's \texttt{*} matches the empty string, and
the two readings cannot both hold: a glob has exactly one wildcard, which either requires
a character or does not. That is not a detail about two patterns. Re-deciding the corpus
with the constant flipped --- a measurement of this paper's procedure under the other
reading, not of what that reading makes expressible, since
Proposition~\ref{prop:decidable}'s search assumes the reading taken --- refuses $599$ of
the $1{,}041$ concrete occurrences where the reading taken refuses $84$. Of the $184$
patterns, $72$ change classification and $515$ occurrences with them, every one of them
toward refusal, because a star that cannot stand for nothing cannot stand for the
\texttt{.*} most of these patterns are built from. It was put to a real account rather than
guessed, and the wildcard does match the empty string, so the figures above stand as
computed.

\paragraph{What the reverse direction costs.}
The practical conclusion is wider than Table~\ref{s:tab:regex-glob} suggests on its own,
and narrower than the refusal count suggests. A Datadog scope reaches $91.9\%$ of what
this corpus's concrete matchers say, but two thirds of that is reached only as a union, and
a union of $41$ branches is a scope somebody has to maintain --- expressible and unpleasant
are different properties, and only the first is what a theorem decides. What is left over
is small and of two kinds. \texttt{2..|5..} is refused by the branch bound, which this
paper chose; raise the bound and it decides, at the price of a scope nobody would write.
\texttt{5[0-9]+} is refused by a theorem, and no bound removes it. That second kind is the
gap this paper set out to distinguish from unfinished work, and in this corner of the
problem there is no unfinished work at all: every refusal is a theorem about globs, or that
theorem together with a bound this paper states.
